\begin{thebibliography}{99}
\bibitem{hmtg26radial} \emph{Rigidité simultanée du lecteur radial ;
composition métrique, mémoire et couplage}. Développements HMT--MD des
25--26 septembre 2026. Démonstrations incorporées dans
\ref{g26:sec:rigidez-radial-es} ; textes intégraux et vérificateurs dans
\path{reproduccion_gravitatoria_20260926/demostraciones} et
\path{reproduccion_gravitatoria_20260926/pruebas}.
\bibitem{hmtg26fase} \emph{Action couplée, phase et mémoire de frontière ;
rigidité et explication structurelle}. Développements HMT--MD des
25--26 septembre 2026. Les démonstrations conservent la phase
relevée, les blocs croisés du poids et les conditions de la
réalisation radiale. La reproduction locale distingue les vérifications
exactes, l'évaluation décimale et les hypothèses de réalisation.
\phantomsection
\addcontentsline{toc}{section}{Références}
\bibitem{mustovicary} B. Musto et J. Vicary, ``Quantum Latin Squares and Unitary Error Bases'', \emph{Quantum Information and Computation} \textbf{16}, 1318--1332 (2016). \href{https://doi.org/10.26421/QIC16.15-16-4}{doi:10.26421/QIC16.15-16-4}.


\bibitem{delascuevas2020} G. De las Cuevas, T. Drescher et T. Netzer, ``Quantum Magic Squares: Dilations and Their Limitations'', \emph{Journal of Mathematical Physics} \textbf{61}, 111704 (2020). \href{https://doi.org/10.1063/5.0022344}{doi:10.1063/5.0022344}.


\bibitem{delascuevas2023} G. De las Cuevas, T. Netzer et I. Valentiner-Branth, ``Magic Squares: Latin, Semiclassical, and Quantum'', \emph{Journal of Mathematical Physics} \textbf{64}, 022201 (2023). \href{https://doi.org/10.1063/5.0127393}{doi:10.1063/5.0127393}.


\bibitem{bennett2003} C. H. Bennett, ``Notes on Landauer's principle,
reversible computation, and Maxwell's Demon'', \emph{Studies in History
and Philosophy of Modern Physics} \textbf{34}, 501--510 (2003).
\href{https://doi.org/10.1016/S1355-2198(03)00039-X}{doi:10.1016/S1355-2198(03)00039-X}.
\bibitem{flm} I. B. Frenkel, J. Lepowsky et A. Meurman, \emph{Vertex Operator Algebras and the Monster}, Pure and Applied Mathematics, vol. 134, Academic Press, 1988.
\bibitem{borcherds} R. E. Borcherds, ``Monstrous Moonshine and Monstrous Lie Superalgebras'', \emph{Inventiones Mathematicae} \textbf{109}, 405--444 (1992). \href{https://doi.org/10.1007/BF01232032}{doi:10.1007/BF01232032}.
\bibitem{conwaysloane} J. H. Conway et N. J. A. Sloane, \emph{Sphere Packings, Lattices and Groups}, 3\textsuperscript{e} éd., Grundlehren der mathematischen Wissenschaften 290, Springer, 1999. \href{https://doi.org/10.1007/978-1-4757-6568-7}{doi:10.1007/978-1-4757-6568-7}.


\end{thebibliography}
